\section{Un tablier tiré des notes de Grothendieck: les pseudo-corps}\label{sec:pseudofields}

Les précédents du \S\ref{sec:precedents} montrent Grothendieck lisant des invariants d'un topos sur deux présentations. Ses notes d'algèbre universelle contiennent aussi un tablier que le registre ne contenait pas encore: le topos des \q{pseudo-corps}, dont il donne trois présentations.

\subsection{Le tablier}

Dans le dossier 161-2 \cite{grothendieck-161-2}, pp.~38--43, Grothendieck énumère des propriétés d'anneaux et, pour chacune, un sous-topos du topos $X_{\mathrm{ann}}\simeq\widehat R$ qui classifie les anneaux, $R$ désignant la catégorie des schémas affines de type fini sur $\mathbb Z$. Sa numérotation change d'une page à l'autre: à la p.~41, l'item~12 est celui des corps, tandis qu'à la p.~42 il renumérote, de sorte que l'item~12 devient celui des anneaux compacts et l'item~13 celui des anneaux compacts réduits. Dans la suite, \q{item~12} et \q{item~13} renvoient à la numérotation de la p.~42, et nous parlons des \q{corps de la p.~41}. L'item~13 est l'\emph{anneau compact réduit}, qu'il nomme \emph{pseudo-corps}: pour tout $f$, il existe un unique idempotent $e$ tel que $ef$ soit inversible dans $eA$ et $(1-e)f=0$, de sorte que $A\simeq A_f\times A/fA$. Ce sont les anneaux commutatifs réguliers au sens de von Neumann. Il donne trois présentations du sous-topos qu'ils définissent:
\[
X_{\mathrm{ps.\,corps}}\;\simeq\;\widetilde R_{\mathrm{ps.\,corps}}\;\simeq\;\widehat{R_{\mathrm{cons}}},
\]
les faisceaux sur $R$ pour la topologie où une famille est couvrante quand l'un de ses membres a une section fibre par fibre, et les préfaisceaux sur la catégorie $R_{\mathrm{cons}}$ obtenue à partir de $R$ en inversant les flèches bijectives à extensions résiduelles triviales (pp.~41--42).

\subsection{Vérification de la présentation par préfaisceaux}

La théorie des pseudo-corps s'axiomatise par les axiomes des anneaux commutatifs et le séquent
\[
\forall a\ \exists!\,b\ \ (a^2b = a\ \wedge\ ab^2 = b):
\]
dans un anneau commutatif, le quasi-inverse $b$ de $a$ est unique, et l'idempotent de la définition de Grothendieck est $e = ab$. La théorie est donc cartésienne, et par suite de type préfaisceau: son topos classifiant est le topos des foncteurs à valeurs dans $\Set$ sur ses modèles de présentation finie. Notons $T$ l'adjoint à gauche de l'inclusion des anneaux réguliers au sens de von Neumann dans les anneaux. Cette inclusion commute aux colimites filtrantes, donc $T$ préserve la présentation finie. Réciproquement, tout pseudo-corps de présentation finie est de la forme $T(A)$, avec $A$ un anneau de présentation finie: dans une présentation finie, on remplace chaque quasi-inverse par une nouvelle variable soumise à ses deux équations. D'après Olivier \cite{olivier}, $\operatorname{Spec}T(A)$ est $\operatorname{Spec}A$ muni de sa topologie constructible, et les anneaux locaux de $T(A)$ sont les corps résiduels de $A$. D'après Pierce \cite{pierce}, un anneau régulier au sens de von Neumann est l'anneau des sections globales du faisceau de ses anneaux locaux sur son spectre, qui est un espace booléen. Ainsi, pour $u:A\to B$ de présentation finie, $T(u)$ est un isomorphisme si et seulement si $\operatorname{Spec}u$ est bijectif et toutes les extensions résiduelles sont triviales: une bijection continue entre espaces booléens est un homéomorphisme, et un morphisme de faisceaux qui est un isomorphisme sur les fibres est un isomorphisme. Les flèches que $T$ inverse sont donc exactement celles de la classe $W$ qu'inverse Grothendieck.

Reste à voir que les pseudo-corps de présentation finie s'obtiennent à partir des anneaux de présentation finie en inversant $W$; voici notre argument. La construction d'Olivier écrit $T(B)$ comme colimite filtrante de $B$-algèbres de présentation finie obtenues en adjoignant des quasi-inverses, un élément $a$ à la fois:
\[
B' = B[y]/(a^2y - a,\ ay^2 - y).
\]
Chaque morphisme $B\to B'$ de cette tour est bijectif sur les spectres, à extensions résiduelles triviales: au-dessus d'un point où $a$ est non nul dans le corps résiduel, les équations imposent $y = 1/a$, et là où $a = 0$, elles imposent $y = 0$. Ces morphismes sont donc dans $W$, de même que leurs composés. Comme $A$ est de présentation finie, un morphisme $T(A)\to T(B)$, c'est-à-dire un morphisme $A\to T(B)$, se factorise par un étage de la tour; il est donc représenté par un toit $A\to B'\leftarrow B$, avec $B\to B'$ dans $W$. La catégorie des pseudo-corps de présentation finie est ainsi la localisée de la catégorie des anneaux de présentation finie par rapport à $W$; son opposée est $R_{\mathrm{cons}}$, et le topos classifiant est $\widehat{R_{\mathrm{cons}}}$, comme l'énonce Grothendieck. Nous n'avons pas vérifié la première présentation, par la topologie $\mathrm{ps.\,corps}$, et rien dans cette section n'est formalisé.

\subsection{Les arches sur ce tablier}

\begin{itemize}
\item \emph{De Morgan:} oui, car les modèles de présentation finie d'une théorie cartésienne ont des sommes amalgamées, et donc s'amalgament (pont~\ref{br:dm}).
\item \emph{Deux valeurs:} non, car l'énoncé $1=0$ définit un sous-objet de l'objet final qui n'est ni $0$ ni $1$.
\item \emph{Connexe:} oui, car $T(\mathbb Z)$ est un modèle de présentation finie initial, de sorte que la catégorie de ces modèles est connexe.
\item \emph{Booléen:} non, car le morphisme $T(\mathbb Z)\to T(\mathbb F_2)$ n'est pas inversible.
\item \emph{Points essentiels:} les pseudo-corps de présentation finie $T(A)$, à isomorphisme près.
\item \emph{Sous-topos:} les théories quotients de la théorie des pseudo-corps, en vertu de la dualité entre théories quotients et sous-topos \cite{caramello-tst}. Plusieurs items de la liste de Grothendieck en font partie: les corps de la p.~41, qu'il identifie aux pseudo-corps intègres et aux pseudo-corps locaux (il marque la dernière égalité \q{à prouver}, et elle est vraie, puisqu'un anneau local régulier au sens de von Neumann est un corps); les corps séparablement clos et algébriquement clos (items~14--15); les anneaux annulés par $n$ et ceux où $n$ est inversible (item~16, p.~43), qu'il marque comme sous-topos ouverts.
\end{itemize}

\subsection{Réponse à une question de la marge}

Pour l'item~12, les anneaux \emph{compacts} (tout $f$ inversible sur un idempotent $e$ et nilpotent sur $1-e$, p.~42), Grothendieck note en marge que la topologie n'est pas cohérente et demande si le topos a assez de points. À la p.~41, il avertit aussi que les sous-topos de sa liste sont déterminés par leurs points, mais que ce n'est pas clair pour leurs intersections, de sorte que vérifier une inclusion entre intersections sur les structures ensemblistes pourrait ne pas suffire. Nous levons ces deux doutes en passant par les axiomes. Chaque item de sa liste s'axiomatise, sur la signature dénombrable des anneaux, par un nombre fini de séquents géométriques. Pour les anneaux compacts, l'un de ces séquents comporte une disjonction dénombrable, puisque $(1-e)f$ est nilpotent d'un certain ordre $n$, sans borne sur $n$; c'est pourquoi la topologie n'est pas cohérente. Le topos classifiant d'une théorie géométrique dénombrable est le topos des faisceaux sur son site syntactique, catégorie dénombrable à limites finies dont la topologie est engendrée par une famille dénombrable de cribles. Or un théorème de Makkai et Reyes \cite[th.~6.2.4]{makkai-reyes} (voir aussi \cite[cor.~3.3.10]{diliberti-rogers}) affirme que le topos des faisceaux sur une catégorie dénombrable à produits fibrés, pour une topologie engendrée par une famille dénombrable de cribles, a assez de points. Le topos des anneaux compacts a donc assez de points; celui des pseudo-corps en a trivialement assez, puisque c'est un topos de préfaisceaux. Dès qu'un sous-topos $E'$ a assez de points, une inclusion $E'\subseteq E''$ se vérifie sur les points. Une intersection d'une famille dénombrable de sous-topos de sa liste classifie la théorie obtenue en réunissant leurs axiomes, qui reste dénombrable; elle a donc assez de points, et les vérifications que Grothendieck fait sur les points sont légitimes pour elle. L'argument ne couvre pas les intersections d'une famille non dénombrable de sous-topos.
